\section*{Arquitectura generadora y conservación de la información}
\addcontentsline{toc}{section}{Arquitectura generadora y conservación de la información}
\label{libnar:troncal}

La Holografía Modular Triádica, abreviada HMT, sitúa el origen de la
construcción en un sistema finito de posiciones, operaciones y transportes.
De su prolongación se obtienen historias compatibles; de las historias,
lecturas numéricas e incidenciales; de estas realizaciones, las estructuras
con las que opera la Mecánica Dimensional. La pregunta que organiza este
artículo concierne a esa continuidad constructiva: qué permanece cuando una
unidad cambia de escala, cuando una operación modifica su representación
y cuando una parte del estado se transfiere a la memoria.

El contenido de la unidad se precisa a lo largo del recorrido. En una
partición aritmética, se conserva la suma normalizada de sus componentes.
En un espacio con producto interno, se conservan normas y correlaciones.
En una representación geométrica, se transportan formas, direcciones e
incidencias. Cada conservación tiene un objeto y una ecuación propios.
Su articulación permite seguir una misma construcción desde el acarreo
de una cifra hasta la recuperación de un estado distribuido en un archivo
de profundidad arbitraria.

La exposición siguiente presenta ese recorrido en su orden de dependencia.
Las definiciones y demostraciones completas aparecen después, dentro del
cuerpo matemático. Aquí interesa que el lector pueda reconocer el papel
de cada operación antes de estudiar su formulación general: cómo se
produce una historia, qué determina una lectura y de qué manera puede
reconstruirse la información que una representación visible ha dejado
de mostrar.

\subsection*{De las posiciones finitas a las historias operatorias}
\label{libnar:app}

La construcción denominada APP, por \emph{All Possible Paths}, parte de
dos ejes de nueve marcas y de dos operaciones sobre sus pares de
posiciones: suma y producto. Las posiciones tienen una identidad que se
conserva al evaluarlas. El par \((i,j)\) admite la lectura aditiva
\(i+j\) y la multiplicativa \(ij\); una ruta registra además cómo se
ha llegado a ese par, qué operación ha actuado y qué orientación ha
seguido el transporte. La relación entre acumulación aditiva y
composición multiplicativa queda así situada en un soporte común.

El paso modular conserva tanto el resto como el cociente. Para una
evaluación positiva \(n\), la marca nonádica \(r\), comprendida entre
uno y nueve, y el cociente \(q\) satisfacen
\[
 n=r+9q.
\]
Esta igualdad permite recuperar la evaluación entera desde su lectura
modular enriquecida. Dos posiciones pueden compartir la misma marca y
conservar cocientes diferentes. Por ejemplo, \((5,5)\) y \((8,2)\)
producen la suma diez; sus productos son veinticinco y dieciséis. Ambos
productos tienen marca siete, con cocientes dos y uno, respectivamente.
La diferencia operatoria queda localizada en el registro que acompaña a
la marca. Las definiciones de la \cref{sec:nucleo} desarrollan esta
conservación conjuntamente con las hojas de operación y las rutas.

La unidad ternaria orientada, denominada TRIT, discrimina tres regímenes
mediante los signos \(-1,0,+1\). Su acción incluye la elección de fase,
la orientación de las prolongaciones y el acarreo. La división ternaria
balanceada escribe un entero como \(n=r+3c\), con
\(r\in\{-1,0,+1\}\); el resto identifica el estado local y el
cociente conserva lo transportado al nivel siguiente. Las realizaciones
geométricas posteriores distinguen los regímenes elíptico, parabólico e
hiperbólico. En la lectura temporal del corpus, el retorno con memoria,
el umbral y la apertura bidireccional poseen acciones específicas sobre
el estado.

El \emph{Topological Phase Kernel}, o núcleo topológico de fase, se
abrevia TPK. Compone la selección trítica, el transporte y la
actualización de memoria. Una transición modifica posiciones, hojas o
fase, y conserva los datos necesarios para relacionar el estado anterior
con el siguiente. Su resultado es un estado enriquecido: contiene
residuos, cocientes, acarreos, orientación, frontera y prefijos de ruta.
El número que se leerá posteriormente procede de esa historia
operatoria. Su genealogía conserva las operaciones que permiten
continuarlo y las relaciones que permiten compararlo con otras
historias.

Esta organización da un significado preciso al retorno. El calendario
de ciento ocho posiciones se distribuye en doce ventanas de nueve. Al
completar una vuelta nonádica retorna la fase correspondiente y avanza
el registro de ventana. El regreso a una posición visible y la
recuperación del estado enriquecido son, por tanto, dos sucesos
diferenciados. La \cref{sec:extension} formaliza el calendario, la
holonomía y el incremento de memoria que acompaña a sus retornos.

\subsection*{La contracción del valor y la permanencia de la incidencia}
\label{libnar:continuo}

Una historia se prolonga mediante prefijos compatibles. Cada prefijo
conserva las operaciones efectuadas y delimita sus continuaciones
admisibles. La construcción conjunta del continuo mantiene en ese mismo
sistema el transporte graduado, los balances de memoria, las incidencias
ternarias, las relaciones de frontera y las composiciones
determinantales. Estas operaciones se enlazan por sus acciones sobre
las mismas historias y por sus compatibilidades con el refinamiento
(\cref{iv:continuo-conjunto}). La correlación entre ellas es parte de
la construcción del objeto.

La lectura arquimediana asocia a los prefijos cilindros de precisión
creciente. Sus intervalos compatibles se contraen y determinan un valor
cuando sus diámetros tienden a cero. La convención semiabierta distingue
los representantes de frontera durante el refinamiento; la información
de acarreo registra las continuaciones que alcanzan un mismo extremo.
En esta lectura aparece la realización numérica del continuo generado.
La secuencia de intervalos expresa cuánto se ha determinado del valor;
el estado enriquecido expresa cómo se ha determinado.

La lectura incidencial conserva otra información: qué elementos forman
un bloque, cómo se orienta una frontera, qué transportes relacionan los
prefijos y qué acciones preservan esas relaciones. A través de los mapas
desarrollados en la \cref{exc:seccion}, esta información se realiza en
matrices de Paley y Hadamard, diseños de Witt y sistemas de Steiner,
códigos de Golay, construcciones reticulares y el retículo de Leech.
La prolongación al álgebra de operadores de vértice \(V^\natural\)
sitúa allí la acción del grupo Monstruo. Cada paso conserva su portador,
su aplicación y sus condiciones de compatibilidad.

La relación entre número y simetría se vuelve así concreta. El valor
arquimediano identifica historias bajo una equivalencia numérica; la
incidencia conserva relaciones entre las historias y sus fronteras.
Ambas lecturas proceden de un mismo estado y cumplen funciones
complementarias. Una magnitud escalar puede permanecer fija mientras
la representación incidencial experimenta una transformación. Seguir
las dos lecturas permite estudiar conjuntamente la estabilización del
valor y la evolución de la estructura que lo realiza.

La unidad también se conserva al distribuirse entre prolongaciones.
Las medidas de los cilindros sucesores suman la medida del cilindro predecesor.
En los espacios de funciones asociados, esta igualdad se convierte en
una isometría de refinamiento: una función del nivel anterior conserva
su norma al representarse sobre las nuevas posiciones. El mismo
principio relaciona partición, medida y transporte. Las demostraciones
de la construcción conjunta especifican la composición y el paso al
límite, de modo que cada nueva profundidad permanece vinculada a las
anteriores.

La realización de una medida añade una operación concreta a esta relación
entre estado y lectura. El sistema y el dispositivo se acoplan, producen
un registro y establecen qué distinciones resultan accesibles. En una
realización isométrica del acoplamiento, la descripción conjunta conserva
las correlaciones; el resultado registrado determina después la
actualización correspondiente. Evaluar un observable, condicionar el
estado a un resultado y preparar nuevamente el registro tienen funciones
diferenciadas. Esta secuencia permite seguir tanto la información que
llega a la lectura como la que permanece en el estado y en el dispositivo.
La memoria interviene en esa continuidad operativa entre una observación
y las siguientes (secciones~\ref{nuclear:3} y~\ref{lib04:archivo}).

\subsection*{La unidad dimensional y el acarreo recuperable}
\label{libnar:unidad}

La extrema y media razón holográfica se estudia aquí a través de las
operaciones que conservan la unidad durante su extensión y su
refinamiento. La posición añadida tiene una identidad propia. Si
\(E_9\) contiene las nueve clases nonádicas, su completación se escribe
\(E_9\sqcup\{\ast\}\). La clase residual cero ya pertenece a
\(E_9\) y tiene el representante positivo nueve; \(\ast\) designa
el estado adicional de frontera. La décima posibilidad conserva así una
función distinta de la reducción modular. En varias coordenadas, contar
cuántas han alcanzado \(\ast\) determina los estratos de la completación
y permite seguir la contribución de cada uno a la unidad.

La completación cúbica ofrece una primera representación
exacta. Al ampliar de nueve a diez las posiciones de cada coordenada,
la unidad cúbica se descompone en interior, estratos de ampliación y
vértice final:
\[
 10^3=9^3+3\cdot9^2+3\cdot9+1
     =729+243+27+1=729+271.
\]
Dividir por mil convierte el recuento en una partición de masa unidad.
Los términos conservan su procedencia dimensional: una, dos o tres
coordenadas ocupan la posición añadida. La expansión
\((9+1)^d\) conserva esta distinción para cada dimensión entera
positiva. El aumento de dimensión y el cambio de escala quedan
relacionados por mapas de coordenadas y sumas de estratos
(\cref{sec:revision-emrh,lib03:completacion}).

El acarreo introduce una transformación más fina. Consideremos la
partición \(73+27=100\). La multiplicación por diez produce
\(730+270=1000\). Transferir una unidad de la primera componente a
la segunda da
\[
 (73,27,1)\longmapsto(10\cdot73-1,10\cdot27+1)
                  =(729,271).
\]
El tercer dato es el acarreo. La suma final sigue siendo mil y la
partición normalizada sigue siendo uno. A la vez, el residuo decimal
de la segunda componente conserva la transferencia: de \(271\) se
recupera el resto uno, después \((271-1)/10=27\), y finalmente
\((729+1)/10=73\). El resultado lleva consigo los datos de la
operación que lo ha producido.

El teorema~\ref{lib03:teo-residuo} demuestra esta inversión para la aplicación
\[
 L_c(x,y)=(Bx-c,By+c),\qquad x+y=M,
\]
con base \(B\), acarreo canónico \(0\leq c<B\) y las condiciones
de positividad declaradas. Tras varias etapas, el registro acumulado
\(C_n\) reúne los acarreos como una palabra en base \(B\), y las
componentes satisfacen
\[
 x_n=B^n x_0-C_n,\qquad y_n=B^n y_0+C_n.
\]
La suma conserva \(B^nM\). Fijadas base y profundidad, el estado
final recupera la palabra completa de acarreos y la partición inicial
(\cref{lib03:cor-palabra}). La profundidad forma parte del dato de
reconstrucción, pues determina también las posiciones iniciales que
contienen ceros.

Las lecturas normalizadas convergen con cotas explícitas para la cola.
Así se enlazan tres niveles de descripción: una operación entera
reversible, una sucesión de particiones de la unidad y un sistema de
cilindros que determina una lectura límite. La elección efectiva de las
continuaciones sigue las reglas del transporte considerado. El ejemplo
\((729,271)\) muestra su primera transferencia; las pruebas de
refinamiento describen también la memoria de todas las transferencias
ulteriores y las coincidencias de frontera entre expansiones
(\cref{lib03:teo-limite,lib03:teo-cilindros}).

\subsection*{Las regiones numéricas y el registro de acoplamiento}
\label{libnar:registro}

En el corpus, las coordenadas \(\pi_{\rm HMT}\),
\(\varphi_{\rm HMT}\) y \(e_{\rm HMT}\) proceden de lectores
regionales sobre la prolongación APP--TRIT--TPK. La región de cierre,
la propagación y la relación estable entre modos poseen operadores
propios. Por ejemplo, el calendario de modos aditivo y multiplicativos
produce la matriz de incidencia
\(F_{\rm av}=\bigl(\begin{smallmatrix}0&1\\1&1\end{smallmatrix}\bigr)\).
Su dirección positiva estable determina la coordenada áurea en la
realización correspondiente. Las pruebas de generación precisan, para
cada profundidad solicitada, cómo se obtiene un prefijo certificado y
cómo se conserva su compatibilidad con los siguientes
(\cref{sec:generacion}).

El mismo estado terminal conserva un registro orientado de sus doce
ventanas. La extracción de los canales firmados, la carga total y la
inversión de Hadamard producen el vector ordenado \(K\). El término
\emph{constante holográfica de acoplamiento aritmético-geométrico}
designa ese registro con sus lectores y con el origen de fase fijado.
Su función consiste en relacionar una codificación aritmética exacta
con las realizaciones de incidencia y geometría que actúan sobre el
mismo objeto (\cref{lib01:definicion}).

En la representación integral canónica se obtiene
\[
 \begin{aligned}
 K=\bigl(&234,543,140,729,659,824,\\
         &621,58,914,794,146,601\bigr),
 \qquad\sum_{j=1}^{12}K_j=6263.
 \end{aligned}
\]
La \cref{lib01:registro} muestra la extracción completa, las
congruencias de sus canales y la divisibilidad que permite invertir
la matriz de Hadamard. El orden de sus componentes conserva la
posición de ventana. Esa información interviene tanto en la lectura
de cifras como en los operadores cíclicos y en las incidencias.

La suma total cumple una función de reconstrucción. Las diferencias
entre posiciones separadas por tres y cuatro pasos determinan
conjuntamente las variaciones relativas del registro. Añadir una misma
cantidad a las doce componentes conserva esas diferencias; la suma fija
precisamente ese nivel común. En \(K\), la carga \(Q(K)=6263\)
restituye la componente uniforme. Diferencias, carga y compatibilidades
enteras permiten recuperar el registro completo
(\cref{prop:k-transversal}). La relación global entre las posiciones y
la organización relativa de sus canales quedan conservadas por lecturas
complementarias.

En base mil, con doce posiciones declaradas, el registro se codifica
por la fracción periódica
\[
 \kappa_{\rm ag}=
 \frac{234543140729659824621058914794146601}{10^{36}-1}.
\]
Cada componente ocupa tres cifras; por ello, la componente cincuenta
y ocho se escribe \(058\) en la expansión. Multiplicar la fracción
por su denominador y efectuar doce divisiones euclídeas recupera el
registro exacto. La periodicidad de esta codificación describe una
lectura fijada del registro finito; la actualización de la historia
enriquecida conserva, por su parte, su propia dinámica de memoria.

La conexión con la coordenada \(\alpha_A\) se expresa en el lector
de clausura por
\[
 \alpha_A=
 \pi_{\rm HMT}+e_{\rm HMT}-\varphi_{\rm HMT}-4-\kappa_{\rm ag}.
\]
Las tres coordenadas regionales y el registro proceden del estado
generado y actúan en el orden constructivo descrito en
\cref{lib01:clausura}. La fórmula hace visible su acoplamiento
aritmético: la extracción de \(K\) conserva una parte de la
estructura que interviene en esa clausura. Las otras realizaciones
de \(K\) conservan sus lectores respectivos. De este modo, una
identidad escalar, una matriz de incidencia y un transporte de fase
pueden estudiarse desde la misma procedencia, con sus acciones
claramente diferenciadas.

\subsection*{Un registro capaz de identificar una transformación}
\label{libnar:identificacion}

La función del registro puede examinarse mediante una operación
concreta. Sea \(S\) el desplazamiento cíclico de sus doce
componentes. Las doce columnas
\(K,SK,\ldots,S^{11}K\) forman una matriz \(O_K\) invertible.
El resultado significa que la órbita cíclica de este registro alcanza
todas las direcciones de su espacio de doce componentes. La prueba
calcula el espectro del Gram \(O_K^*O_K\) y da cotas positivas
exactas para sus valores propios (\cref{lib01:marco-ciclico}).

Supongamos que un operador \(H\) respeta el desplazamiento cíclico,
es decir, \(HS=SH\). Al observar la respuesta vectorial \(y=HK\),
sus desplazamientos proporcionan las respuestas a todas las columnas
de \(O_K\). Se recupera entonces el operador completo mediante
\[
 H=O_yO_K^{-1},\qquad
 O_y=(y,Sy,\ldots,S^{11}y).
\]
La observación contiene doce componentes, y la hipótesis de
conmutación identifica la clase de transformaciones recuperables.
Dentro de esa clase, una sola respuesta vectorial determina todos
los coeficientes del filtro circular. Por ejemplo, la combinación
\(2I+3S-S^5\) queda identificada desde su acción sobre \(K\),
con las cotas de error del teorema~\ref{lib01:operadores}.

Este hecho ofrece una aplicación operativa de la estructura del
registro. Su orden y su espectro permiten usarlo como estado de
prueba para identificar transformaciones compatibles. Si el estado
y sus respuestas se transportan mediante una isometría completa de
memoria, el Gram de la órbita se conserva y la identificación sigue
siendo recuperable por el adjunto. La estabilidad procede de valores
propios calculados y de operadores inversores explícitos. El mismo
registro que participa en la clausura aritmética adquiere así una
función precisa en el reconocimiento de transportes.

La normalización también tiene una información asociada. El factor
unitario de \(O_K\) conserva orientación y ángulos; su Gram
conserva las amplitudes necesarias para reconstruir la matriz
original. El par formado por ambos factores recupera \(O_K\)
(\cref{lib01:polar}). Esta separación permite elegir una
representación normalizada y mantener, al mismo tiempo, la escala
exacta del registro.

\subsection*{Dos transportes y una conservación bilateral}
\label{libnar:bilateral}

El siguiente paso relaciona el refinamiento con una ley de
conservación de correlaciones. Sean \(\mathcal B\) y
\(\mathcal C\) dos isometrías que transportan el mismo espacio de
estados a un espacio de salida común. Para \(a>0\), se forman los
dos operadores
\[
 Q_-=a\mathcal B-\mathcal C,\qquad
 Q_+=(a+1)\mathcal B+\mathcal C.
\]
Las dos combinaciones reúnen una contribución común y una
contribución relativa con signos diferentes. Al sumar sus normas
con los pesos correspondientes, los términos cruzados se cancelan:
\[
 (a+1)\|Q_-v\|^2+a\|Q_+v\|^2
   =(2a+1)(a^2+a+1)\|v\|^2.
\]
El teorema~\ref{lib04:teo-balance} demuestra esta igualdad para cualquier
dimensión de los espacios y para los transportes isométricos
declarados. La misma cancelación conserva productos internos y
admite su formulación para formas preservadas por los transportes.

En el caso \(a=8\), el factor \(a^2+a+1\) vale setenta y tres.
La factorización de la isometría completa distribuye primero la
norma entre dos componentes con pesos
\[
 \frac{72}{73}\qquad\text{y}\qquad\frac{1}{73}.
\]
Después actúa una mezcla ortogonal que produce los canales
normalizados de \(Q_-\) y \(Q_+\). Los pesos tienen así un
origen operatorio preciso: \(72=8\cdot9\) y
\(73=8\cdot9+1\). Su función se conserva al cambiar la dimensión
o los transportes, dentro de las hipótesis de la ley bilateral
(\cref{lib04:factorizacion}).

La recuperación se ve directamente en las combinaciones de salida.
Si \(x=Q_-v\) e \(y=Q_+v\), entonces
\[
 x+y=(2a+1)\mathcal Bv,\qquad
 ay-(a+1)x=(2a+1)\mathcal Cv.
\]
Los dos transportes del estado inicial permanecen determinados por
el par de canales. El adjunto de cualquiera de las isometrías
recupera después el estado correspondiente. Con la normalización
completa, el operador \(W_a\) satisface \(W_a^*W_a=I\):
conserva la norma y ofrece una inversa exacta sobre su imagen.

Una lectura particular puede variar mientras se conserva toda esa
información. Cuando \(\mathcal C=\mathcal B R\), con
\(\mathcal B\) y \(R\) unitarios, este último como transporte relativo, la
lectura diagonal de un observable \(G\) en ambos canales produce
\[
 \mathcal T(G)=\frac{72G+R^*GR}{73}
 \qquad(a=8).
\]
Los observables compatibles con \(R\) permanecen fijos; los demás
registran la diferencia de sus dos acciones. Esta ecuación distingue
la conservación del estado completo y la respuesta de un lector
determinado (\cref{lib04:teo-lector}).

\subsection*{El mismo acarreo, ahora como estado transportable}
\label{libnar:ejemplo}

La relación anterior se aplica materialmente al acarreo. A cada
etiqueta admisible \((x,y,c)\) se asocia un vector de una base
ortonormal. La transformación entera \(L_c\) induce entonces el
operador \(J\) que lleva esa base a las etiquetas de salida
\((Bx-c,By+c)\). Su inversa recupera el acarreo mediante el
residuo. La biyección de etiquetas convierte a \(J\) en una
transformación unitaria (\cref{lib05:prop-unitaria}). Los valores
enteros designan las etiquetas; las normas corresponden a las
amplitudes de los vectores que las representan.

Para base diez y suma inicial cien, consideremos las dos etiquetas
\(d_1=(73,27,1)\) y \(d_2=(73,27,2)\). Sus imágenes son
\(e_1=(729,271)\) y \(e_2=(728,272)\). Fijemos el transporte
\(R\) que intercambia esas dos etiquetas y deja fijas las demás.
La composición usa \(\mathcal B=J\) y \(\mathcal C=JR\).
Esta elección proporciona un ejemplo explícito dentro del teorema
bilateral, con regla de transporte especificada antes de evaluarlo.

Sobre el estado \(|d_1\rangle\), los canales para \(a=8\)
resultan
\[
 Q_-|d_1\rangle=8|e_1\rangle-|e_2\rangle,
 \qquad
 Q_+|d_1\rangle=9|e_1\rangle+|e_2\rangle.
\]
Sus normas cuadradas son sesenta y cinco y ochenta y dos. El balance
exacto es \(9\cdot65+8\cdot82=1241=17\cdot73\), mientras
que la suma de ambos canales recupera \(17|e_1\rangle\).
Aplicar la inversa de etiquetas recupera la partición \((73,27)\)
y el acarreo uno.

La lectura fraccionaria de la primera componente ofrece además una
predicción exacta para este transporte declarado. Al evaluarla
diagonalmente en ambos canales normalizados se obtiene
\[
 \frac{72\cdot729+728}{73\cdot1000}
   =\frac{729}{1000}-\frac{1}{73000}.
\]
El lector diagonal de acarreo da \((72\cdot1+2)/73=74/73\).
El observable transportado desde la entrada, en cambio, conserva
el valor propio uno del acarreo original. La diferencia entre ambos
resultados permite identificar qué operación de lectura se ha
realizado. La información del estado permanece recuperable y la
lectura diagonal cuantifica cómo interviene el transporte relativo
(\cref{lib05:ejemplo}).

El ejemplo reúne las dos conservaciones del artículo. La suma
aritmética se conserva en las etiquetas refinadas y la norma se
conserva en sus representaciones unitarias. La conexión entre ambas
procede de la biyección y de su levantamiento a bases ortonormales.
Esta composición permite calcular lectores, invertir el transporte
y controlar el error manteniendo identificados los objetos sobre
los que actúa cada ley.

\subsection*{Memoria completa, simetrías y profundidad arbitraria}
\label{libnar:memoria}

La conservación bilateral se prolonga etapa a etapa. En cada nivel,
un canal continúa el estado y el otro registra un complemento de
memoria. Si \(R_Nv\) es el estado que alcanza el nivel \(N\) y
\(m_0,\ldots,m_{N-1}\) son los complementos registrados, el
balance acumulado es
\[
 \|v\|^2=\|R_Nv\|^2+\sum_{n=0}^{N-1}\|m_n\|^2.
\]
La igualdad identifica exactamente dónde permanece la norma
inicial. La recuperación finita utiliza el estado continuado y
todos los complementos disponibles; sus correlaciones se conservan
por la misma isometría conjunta.

Para la política de archivo del teorema~\ref{lib04:teo-archivo}, la
norma del canal continuado posee una cota geométrica uniforme.
En el caso \(a=8\), su norma cuadrada está acotada en cada etapa
por \(729/1241\). Al aumentar la profundidad, la componente
continuada tiende a cero y el archivo infinito conserva toda la
norma. Si \(A_n\) y \(D_n\) designan los bloques de
continuación y memoria, respectivamente, el estado se reconstruye
por la serie convergente
\[
 v=\sum_{n\geq0}R_n^*D_n^*m_n.
\]
El operador inversor tiene norma uno. Una perturbación de norma
\(\varepsilon\) en el archivo completo produce un error de
reconstrucción de norma a lo sumo \(\varepsilon\). Las
truncaciones finitas tienen sus propias cotas y, cuando se emplea
sólo su memoria, la inversa finita especificada en el mismo
desarrollo.

La incidencia del registro \(K\), su órbita cíclica y sus
direcciones distinguidas pueden transportarse mediante estas
isometrías. Se conservan los productos internos que determinan
ángulos y matrices de Gram; los proyectores y lectores se
transportan con sus dominios; la acción del adjunto recupera la
representación inicial. Las potencias exteriores transportan la
orientación y conservan los determinantes de Gram, que miden los
cuadrados de los volúmenes correspondientes. La dimensión del
portador puede ampliarse y la información geométrica de la
imagen permanece especificada por la isometría.

La relación con la conservación de Noether se formula en un nivel
adicional bien determinado. Para la acción discreta auxiliar
construida en la \cref{lib01:carga-noether}, una colección de
generadores que entrelaza los transportes produce una carga
conservada de la forma \(p_n^*A_nq_n\). Aquí intervienen una
acción variacional, sus ecuaciones y la simetría que las preserva.
El balance de norma y la carga variacional se articulan mediante
los transportes comunes, cada uno con sus hipótesis y con la
función que desempeña en la demostración.

El recorrido establece una relación operativa entre unidad,
número, memoria y geometría. La unidad aritmética admite
refinamientos reversibles; el estado enriquecido conserva el
registro de las operaciones; las realizaciones numérica e
incidencial extraen propiedades correlacionadas; el acoplamiento
bilateral mantiene información y correlaciones durante su
redistribución. La extrema y media razón holográfica adquiere así
un desarrollo demostrativo en el que partición, extensión,
transporte y recuperación pueden seguirse mediante aplicaciones
explícitas. Las páginas siguientes presentan esas aplicaciones,
sus pruebas y sus ejemplos completos.
